% book08.tex --- Book VIII of Euclid's Elements: Continued Proportions.
%
% All 27 propositions encoded. Book VIII studies geometric progressions
% of integers, especially their reduction to a least sequence in the
% given ratio.  It is the integer counterpart to Book V's continuous
% theory and feeds directly into Book IX's number-theoretic applications.
%
% Wording follows Heath (1908).

\section{Book VIII --- Continued Proportions}
\label{sec:book-VIII}

\begin{claim}[Proposition VIII.1: Least continued proportion has coprime extremes]
\label{prop:VIII.1}
If there be as many numbers as we please in continued proportion, and
the extremes of them be prime to one another, the numbers are the
least of those which have the same ratio with them.
\end{claim}
\begin{evidence}[Proof of VIII.1]
\label{ev:VIII.1}
Suppose a smaller set $b_1, \dots, b_n$ in the same ratio existed.
By ex aequali (VII.14) the ratio of extremes $b_1 : b_n$ equals $a_1
: a_n$, so by VII.21 (least in a ratio are coprime) the original
$a_1$, $a_n$ would not be coprime --- contradiction.
\dependson{VIII.1}{VII.14}
\dependson{VIII.1}{VII.20}
\dependson{VIII.1}{VII.21}
\end{evidence}

\begin{claim}[Proposition VIII.2: Construct a continued proportion with a given common ratio]
\label{prop:VIII.2}
To find numbers in continued proportion, as many as may be prescribed,
and the least that are in a given ratio.
\end{claim}
\begin{evidence}[Proof of VIII.2]
\label{ev:VIII.2}
For $a : b$ with $\gcd(a, b) = 1$, the sequence $a^{n-1}, a^{n-2}
b, \dots, b^{n-1}$ is in continued proportion with ratio $a : b$; by
VII.27 the extremes are coprime, so by VIII.1 the sequence is least
in its ratio.
\dependson{VIII.2}{VII.27}
\dependson{VIII.2}{VIII.1}
\end{evidence}

\begin{claim}[Proposition VIII.3: Least continued proportion has coprime extremes (converse setup)]
\label{prop:VIII.3}
If as many numbers as we please in continued proportion be the least
of those which have the same ratio with them, the extremes of them
are prime to one another.
\end{claim}
\begin{evidence}[Proof of VIII.3]
\label{ev:VIII.3}
Suppose the extremes shared a common divisor $d > 1$.  By VII.20
each term would be divisible by some power of $d$, producing a
smaller sequence in the same ratio --- contradicting minimality.
\dependson{VIII.3}{VII.20}
\dependson{VIII.3}{VII.21}
\dependson{VIII.3}{VIII.1}
\end{evidence}

\begin{claim}[Proposition VIII.4: Common ratio across multiple chains]
\label{prop:VIII.4}
Given as many ratios as we please in least numbers, to find numbers
in continued proportion which are the least in the given ratios.
\end{claim}
\begin{evidence}[Proof of VIII.4]
\label{ev:VIII.4}
Reduce each ratio to lowest terms by VII.33.  Compound them by
multiplying numerators and denominators across; the resulting
sequence is in continued proportion with the prescribed ratios.
\dependson{VIII.4}{VII.33}
\dependson{VIII.4}{VIII.2}
\end{evidence}

\begin{claim}[Proposition VIII.5: Plane numbers have a compound ratio of their sides]
\label{prop:VIII.5}
Plane numbers have to one another the ratio compounded of the ratios
of their sides.
\end{claim}
\begin{evidence}[Proof of VIII.5]
\label{ev:VIII.5}
For plane numbers $ab$ and $cd$: $ab : cd = (a : c) \cdot (b : d)$
in the language of compound ratios.  Verified by direct computation
using VII.17 / VII.18.
\dependson{VIII.5}{VII.17}
\dependson{VIII.5}{VII.18}
\dependson{VIII.5}{def:VII.16}
\end{evidence}

\begin{claim}[Proposition VIII.6: First measures last iff first measures second]
\label{prop:VIII.6}
If there be as many numbers as we please in continued proportion, and
the first do not measure the second, neither will any other measure
any other.
\end{claim}
\begin{evidence}[Proof of VIII.6]
\label{ev:VIII.6}
Contrapositive of VIII.7: divisibility propagates through the
sequence, so failure at the first step prevents any later
divisibility relation.
\dependson{VIII.6}{VII.20}
\dependson{VIII.6}{VIII.1}
\end{evidence}

\begin{claim}[Proposition VIII.7: First measures last implies first measures second]
\label{prop:VIII.7}
If there be as many numbers as we please in continued proportion, and
the first measure the last, it will measure the second also.
\end{claim}
\begin{evidence}[Proof of VIII.7]
\label{ev:VIII.7}
If $a_1 \mid a_n$, reduce $a_1, \dots, a_n$ to lowest terms (VIII.3);
since the lowest extremes are coprime but $a_1$ divides $a_n$, the
ratio $a_1 : a_n$ must be 1:1 in lowest terms, forcing $a_1 \mid
a_2$ via VII.20.
\dependson{VIII.7}{VII.20}
\dependson{VIII.7}{VIII.3}
\end{evidence}

\begin{claim}[Proposition VIII.8: Intermediate numbers in a proportion]
\label{prop:VIII.8}
If between two numbers there fall numbers in continued proportion
with them, then, however many numbers fall between them in continued
proportion, so many will also fall in continued proportion between
the numbers which have the same ratio with the original numbers.
\end{claim}
\begin{evidence}[Proof of VIII.8]
\label{ev:VIII.8}
The number of geometric means between two numbers depends only on
their ratio; scaling by a common factor changes the magnitudes but
not the ratio, so the same number of means fall between the scaled
pair.
\dependson{VIII.8}{VII.13}
\dependson{VIII.8}{VIII.2}
\end{evidence}

\begin{claim}[Proposition VIII.9: Coprimality between unit and a sequence]
\label{prop:VIII.9}
If two numbers be prime to one another, and numbers fall between them
in continued proportion, then, however many numbers fall between them
in continued proportion, so many will also fall in continued
proportion between each of them and a unit.
\end{claim}
\begin{evidence}[Proof of VIII.9]
\label{ev:VIII.9}
Coprime extremes correspond to a least continued proportion (VIII.1);
the unit extends the proportion at both ends, and VIII.2 gives the
matching extension on each side.
\dependson{VIII.9}{VIII.1}
\dependson{VIII.9}{VIII.2}
\end{evidence}

\begin{claim}[Proposition VIII.10: Counting means between unit and number]
\label{prop:VIII.10}
If numbers fall between each of two numbers and a unit in continued
proportion, however many numbers fall between each of them and a
unit in continued proportion, so many also will fall between them in
continued proportion.
\end{claim}
\begin{evidence}[Proof of VIII.10]
\label{ev:VIII.10}
Concatenate two unit-anchored continued proportions; VII.14 (ex
aequali) confirms the joined sequence remains in continued
proportion.
\dependson{VIII.10}{VII.14}
\dependson{VIII.10}{VIII.9}
\end{evidence}

\begin{claim}[Proposition VIII.11: Squares have one mean proportional]
\label{prop:VIII.11}
Between two square numbers there is one mean proportional number, and
the square has to the square the ratio duplicate of that which the
side has to the side.
\end{claim}
\begin{evidence}[Proof of VIII.11]
\label{ev:VIII.11}
For squares $a^2$, $b^2$: the mean proportional is $ab$ (since $a^2 :
ab = ab : b^2 = a : b$), and $a^2 : b^2$ is the duplicate of $a : b$.
\dependson{VIII.11}{VII.17}
\dependson{VIII.11}{VII.18}
\dependson{VIII.11}{def:VII.18}
\end{evidence}

\begin{claim}[Proposition VIII.12: Cubes have two mean proportionals]
\label{prop:VIII.12}
Between two cube numbers there are two mean proportional numbers,
and the cube has to the cube the ratio triplicate of that which the
side has to the side.
\end{claim}
\begin{evidence}[Proof of VIII.12]
\label{ev:VIII.12}
For cubes $a^3$, $b^3$: the two means are $a^2 b$ and $a b^2$, and
$a^3 : b^3$ is the triplicate of $a : b$.
\dependson{VIII.12}{VII.17}
\dependson{VIII.12}{VII.18}
\dependson{VIII.12}{def:VII.19}
\dependson{VIII.12}{def:V.10}
\end{evidence}

\begin{claim}[Proposition VIII.13: Powers of a continued proportion are in continued proportion]
\label{prop:VIII.13}
If there be as many numbers as we please in continued proportion, and
each by multiplying itself make some number, the products will be
proportional; and if the original numbers by multiplying the products
make certain numbers, the latter will also be proportional.
\end{claim}
\begin{evidence}[Proof of VIII.13]
\label{ev:VIII.13}
Squares (and cubes) of terms in continued proportion are themselves
in continued proportion, by VII.27 and VIII.2.
\dependson{VIII.13}{VII.27}
\dependson{VIII.13}{VIII.2}
\end{evidence}

\begin{claim}[Proposition VIII.14: Square measures square iff side measures side]
\label{prop:VIII.14}
If a square measure a square, the side will also measure the side;
and if the side measure the side, the square will also measure the
square.
\end{claim}
\begin{evidence}[Proof of VIII.14]
\label{ev:VIII.14}
$a^2 \mid b^2 \iff a \mid b$, by Euclid's lemma (VII.30) applied
prime-by-prime.
\dependson{VIII.14}{VII.30}
\dependson{VIII.14}{VIII.11}
\end{evidence}

\begin{claim}[Proposition VIII.15: Cube measures cube iff side measures side]
\label{prop:VIII.15}
If a cube number measure a cube number, the side will also measure
the side; and if the side measure the side, the cube will also
measure the cube.
\end{claim}
\begin{evidence}[Proof of VIII.15]
\label{ev:VIII.15}
Same prime-by-prime argument as VIII.14 for cubes.
\dependson{VIII.15}{VIII.12}
\dependson{VIII.15}{VIII.14}
\end{evidence}

\begin{claim}[Proposition VIII.16: Squares non-measuring]
\label{prop:VIII.16}
If a square measure not a square, neither will the side measure the
side; and if the side measure not the side, neither will the square
measure the square.
\end{claim}
\begin{evidence}[Proof of VIII.16]
\label{ev:VIII.16}
Contrapositive of VIII.14.
\dependson{VIII.16}{VIII.14}
\end{evidence}

\begin{claim}[Proposition VIII.17: Cubes non-measuring]
\label{prop:VIII.17}
If a cube number measure not a cube number, neither will the side
measure the side; and if the side measure not the side, neither will
the cube measure the cube.
\end{claim}
\begin{evidence}[Proof of VIII.17]
\label{ev:VIII.17}
Contrapositive of VIII.15.
\dependson{VIII.17}{VIII.15}
\end{evidence}

\begin{claim}[Proposition VIII.18: Mean proportional between similar plane numbers]
\label{prop:VIII.18}
Between two similar plane numbers there is one mean proportional
number, and the plane number has to the plane number the ratio
duplicate of that which the corresponding side has to the
corresponding side.
\end{claim}
\begin{evidence}[Proof of VIII.18]
\label{ev:VIII.18}
For similar plane numbers $ab$ and $cd$ with $a : b = c : d$, the
mean proportional is the geometric mean of $ab$ and $cd$, which by
VII.19 / VIII.2 equals $ad$ (or $bc$, equal by VII.19).
\dependson{VIII.18}{VII.19}
\dependson{VIII.18}{VIII.5}
\dependson{VIII.18}{def:VII.21}
\end{evidence}

\begin{claim}[Proposition VIII.19: Mean proportionals between similar solid numbers]
\label{prop:VIII.19}
Between two similar solid numbers there fall two mean proportional
numbers, and the solid number has to the solid number the ratio
triplicate of that which the corresponding side has to the
corresponding side.
\end{claim}
\begin{evidence}[Proof of VIII.19]
\label{ev:VIII.19}
For similar solid numbers $abc$ and $def$: the two means are $abf$
and $aef$ (or symmetric variants); together they give a continued
proportion in triplicate ratio.
\dependson{VIII.19}{VIII.12}
\dependson{VIII.19}{VIII.18}
\end{evidence}

\begin{claim}[Proposition VIII.20: Mean proportional characterises similar plane numbers]
\label{prop:VIII.20}
If one mean proportional number fall between two numbers, the numbers
will be similar plane numbers.
\end{claim}
\begin{evidence}[Proof of VIII.20]
\label{ev:VIII.20}
Converse of VIII.18: if $a : m = m : b$ then $a$ and $b$ admit
factorisations as similar plane numbers via VII.19.
\dependson{VIII.20}{VII.19}
\dependson{VIII.20}{VIII.18}
\end{evidence}

\begin{claim}[Proposition VIII.21: Two mean proportionals characterise similar solid numbers]
\label{prop:VIII.21}
If two mean proportional numbers fall between two numbers, the
numbers are similar solid numbers.
\end{claim}
\begin{evidence}[Proof of VIII.21]
\label{ev:VIII.21}
Converse of VIII.19.
\dependson{VIII.21}{VIII.19}
\dependson{VIII.21}{VIII.20}
\end{evidence}

\begin{claim}[Proposition VIII.22: Squares in continued proportion]
\label{prop:VIII.22}
If three numbers be in continued proportion, and the first be square,
the third will also be square.
\end{claim}
\begin{evidence}[Proof of VIII.22]
\label{ev:VIII.22}
If $a^2 : m = m : c$, then $m^2 = a^2 c$ so $c = (m/a)^2$, hence $c$
is square.  VII.19 ensures the division produces an integer.
\dependson{VIII.22}{VII.19}
\dependson{VIII.22}{VIII.11}
\end{evidence}

\begin{claim}[Proposition VIII.23: Cubes in continued proportion]
\label{prop:VIII.23}
If four numbers be in continued proportion, and the first be cube,
the fourth will also be cube.
\end{claim}
\begin{evidence}[Proof of VIII.23]
\label{ev:VIII.23}
Same scheme as VIII.22 with two means; the fourth term is the cube
of the ratio's denominator scaled appropriately.
\dependson{VIII.23}{VIII.12}
\dependson{VIII.23}{VIII.22}
\end{evidence}

\begin{claim}[Proposition VIII.24: Square ratio implies square ratio]
\label{prop:VIII.24}
If two numbers have to one another the ratio which a square number
has to a square number, and the first be square, the second will
also be square.
\end{claim}
\begin{evidence}[Proof of VIII.24]
\label{ev:VIII.24}
By VIII.11 the ratio of squares has a mean proportional; transferring
that mean to $a^2 : b$ forces $b$ to be square by VIII.22.
\dependson{VIII.24}{VIII.11}
\dependson{VIII.24}{VIII.22}
\end{evidence}

\begin{claim}[Proposition VIII.25: Cube ratio implies cube ratio]
\label{prop:VIII.25}
If two numbers have to one another the ratio which a cube number has
to a cube number, and the first be cube, the second will also be
cube.
\end{claim}
\begin{evidence}[Proof of VIII.25]
\label{ev:VIII.25}
Same argument as VIII.24 for cubes via VIII.12 and VIII.23.
\dependson{VIII.25}{VIII.12}
\dependson{VIII.25}{VIII.23}
\end{evidence}

\begin{claim}[Proposition VIII.26: Similar plane numbers have square ratio]
\label{prop:VIII.26}
Similar plane numbers have to one another the ratio which a square
number has to a square number.
\end{claim}
\begin{evidence}[Proof of VIII.26]
\label{ev:VIII.26}
By VIII.18 similar plane numbers admit a mean proportional, and the
ratio (squares of corresponding sides) is a square-to-square ratio.
\dependson{VIII.26}{VIII.18}
\dependson{VIII.26}{def:VII.21}
\end{evidence}

\begin{claim}[Proposition VIII.27: Similar solid numbers have cube ratio]
\label{prop:VIII.27}
Similar solid numbers have to one another the ratio which a cube
number has to a cube number.
\end{claim}
\begin{evidence}[Proof of VIII.27]
\label{ev:VIII.27}
By VIII.19 similar solid numbers admit two mean proportionals, and
the ratio is in the triplicate (cube-to-cube) ratio of corresponding
sides.
\dependson{VIII.27}{VIII.19}
\dependson{VIII.27}{def:VII.21}
\end{evidence}
